\begin{abstract}
%
%
Modern transformer-based Large Language Models (LLMs) are constructed with a series of decoder blocks.
%
Each block comprises three key components: (1) QKV generation, (2) multi-head attention, and (3) feed-forward networks.
%
In batched processing, QKV generation and feed-forward networks involve compute-intensive matrix-matrix multiplications (GEMM), while multi-head attention requires bandwidth-heavy matrix-vector multiplications (GEMV).
%
Machine learning accelerators like TPUs or NPUs are proficient in handling GEMM but are less efficient for GEMV computations.
%
Conversely, Processing-in-Memory (PIM) technology is tailored for efficient GEMV computation, while it lacks the computational power to handle GEMM effectively.


Inspired by this insight, we propose \neupims, a heterogeneous acceleration system that jointly exploits a conventional GEMM-focused NPU and GEMV-optimized PIM devices.
%
The main challenge in efficiently integrating NPU and PIM lies in enabling concurrent operations on both platforms, each addressing a specific kernel type.
%
First, existing PIMs typically operate in a ``blocked'' mode, allowing only either NPU or PIM to be active at any given time.
%
Second, the inherent dependencies between GEMM and GEMV in LLMs restrict their parallel processing.
%
To tackle these challenges, \neupims is equipped with \emph{dual row buffers} in each bank, facilitating the simultaneous management of memory read/write operations and PIM commands.
%
Further, \neupims employs a runtime \emph{sub-batch interleaving} technique to maximize concurrent execution, leveraging batch parallelism to allow two independent sub-batches to be pipelined within a single \neupims device.
%
Our evaluation demonstrates that compared to GPU-only, NPU-only, and a na\"ive NPU+PIM integrated acceleration approaches, \neupims achieves 3$\times$, 2.4$\times$ and 1.6$\times$ throughput improvement, respectively.

\end{abstract}